\documentclass[10pt]{amsart}
\usepackage[a4paper,margin=22mm]{geometry}
\usepackage[no-math]{fontspec}
\setmainfont{Latin Modern Roman}[SmallCapsFont={lmromancaps10-regular.otf}]
\usepackage{amsmath,amssymb,amsthm,mathtools,bbm,hyperref,graphicx,etoolbox}
\emergencystretch=3em
\allowdisplaybreaks
\newtheorem{theo}{Theorem}[section]
\newtheorem{lemm}[theo]{Lemma}
\newtheorem{prop}[theo]{Proposition}
\newtheorem{coro}[theo]{Corollary}
\theoremstyle{definition}
\newtheorem{rema}[theo]{Remark}
\numberwithin{equation}{section}
\newcommand{\N}{\mathbb{N}}
\newcommand{\R}{\mathbb{R}}
\newcommand{\bbB}{\mathbb{B}}
\newcommand{\bbP}{\mathbb{P}}
\newcommand{\bbE}{\mathbb{E}}
\newcommand{\D}{\mathbb{D}}
\renewcommand{\S}{\mathbb{S}}
\renewcommand{\P}{\mathbb{P}}
\renewcommand{\H}{\mathbb{H}}

\DeclareMathOperator{\E}{\mathbb{E}}

\newcommand{\h}{\mathrm{h}}
\newcommand{\dd}{\mathrm{d}}
\DeclareMathOperator{\Hull}{Hull}
\DeclareMathOperator{\nee}{ne}
\DeclareMathOperator{\Ree}{Re}

\newcommand{\calS}{\mathcal{S}}
\newcommand{\calC}{\mathcal{C}}
\newcommand{\calD}{\mathcal{D}}
\newcommand{\calE}{\mathcal{E}}
\newcommand{\calG}{\mathcal{G}}
\newcommand{\calN}{\mathcal{N}}

\newcommand{\bfs}{\mathbf{s}}
\newcommand{\bfe}{\mathbf{e}}

\newcommand{\htec}{\mathbf{H}_{\mathrm{tech}}}
\newcommand{\ind}[1]{\mathbbm{1}_{\left\{#1\right\}}}

\renewcommand{\bar}[1]{\overline{#1}}
\renewcommand{\tilde}[1]{\widetilde{#1}}
\renewcommand{\epsilon}{\varepsilon}




%%%%%%%%%%%%%%%%%%%%%%%%%%%%%%%%%%%%%%%%%%%%%%%%%%%%%%%%%%%%%%%%%%%%%%%%%%%%%%%%%%%%%%%%%%%%%%%%

\graphicspath{{./figures/}}

\newcommand*{\mk}{\mkern -1mu}
\newcommand*{\Mk}{\mkern -2mu}
\newcommand*{\mK}{\mkern 1mu}
\newcommand*{\MK}{\mkern 2mu}

\hypersetup{urlcolor=purple, linkcolor=blue, citecolor=red}

\newcommand*{\relabel}{\renewcommand{\labelenumi}{(\theenumi)}}
\newcommand*{\romanenumi}{\renewcommand*{\theenumi}{\roman{enumi}}\relabel}
\newcommand*{\Romanenumi}{\renewcommand*{\theenumi}{\Roman{enumi}}\relabel}
\newcommand*{\alphenumi}{\renewcommand*{\theenumi}{\alph{enumi}}\relabel}
\newcommand*{\Alphenumi}{\renewcommand*{\theenumi}{\Alph{enumi}}\relabel}
\let\oldtilde\tilde
\renewcommand*{\tilde}[1]{\mathchoice{\widetilde{#1}}{\widetilde{#1}}{\oldtilde{#1}}{\oldtilde{#1}}}
\let\oldhat\hat
\renewcommand*{\hat}[1]{\mathchoice{\widehat{#1}}{\widehat{#1}}{\oldhat{#1}}{\oldhat{#1}}}
\let\oldforall\forall
\renewcommand*{\forall}{\mathrel{\oldforall}}
%%%%%%%%%%%%%%%%%%%%%%%%%%%%%%%%%%%%%%%%%%%%%%%%%%%%%%%%%%%%%%%%%%%%%%%%%%%%%%%%%%%%%%%%%%%%%%%%
\title[Sharp inverse-integral subordinator asymptotics]{Sharp stretched-exponential Laplace asymptotics for inverse-integral subordinators under regular variation and strong non-lattice hypotheses}
\author{B\'en\'edicte Haas}
\author{Bastien Mallein}
\date{Annales Henri Lebesgue8(2025),219–253; DOI10.5802/ahl.234}
\hypersetup{hidelinks}
\AtBeginDocument{\let\originalRef\ref\renewcommand{\ref}[1]{\ifstrequal{#1}{thm:asymp_distribution}{\href{https://doi.org/10.5802/ahl.234}{1.7 (source)}}{\ifstrequal{#1}{sec:proofscaling}{\href{https://doi.org/10.5802/ahl.234}{3.3 (source)}}{\originalRef{#1}}}}}
\begin{document}
\maketitle
\setcounter{section}{1}
\setcounter{theo}{1}
\setcounter{equation}{6}
Motivated by the previous discussion, we consider now a generic subordinator $\xi$, with Laplace exponent
\[
\phi(q)=-\log \E\left[e^{-q \xi_1}\right] = \kappa + c q + \int_0^\infty (1 - e^{-qx}) \pi(\dd x)
\]
where $\kappa \geq 0$, $c \geq 0$ and $\pi$ is a measure on $(0,\infty)$ such that $\int_0^{\infty} (1 \wedge x) \pi(\dd x) < \infty$. We refer to $\pi$ as the Lévy measure of $\xi$, $c$ its drift and $\kappa$ its death rate. We define then the time change $\rho$ by
\[
\rho(t) := \inf \left\{ r \geq 0 : \int_0^r \xi_u \dd u \geq t \right\}, \quad t \geq 0,
\]
with the convention $\inf{\emptyset}=\infty$. Our first main result gives the following exact expression for the Laplace transform of $\xi_{\rho(t)}$, for any $t\geq 0$.
\begin{theo}\label{thm:lapTransform}
Let $\Phi(q)=\int_0^q \phi(s) \dd s$, $q \geq 0$ and $\Phi^{-1}: [0,\infty) \mapsto [0,\infty)$ be the inverse of $\Phi$. Then the function $\Phi^{-1}$ is the Laplace exponent of a subordinator, with a Lévy measure that we denote $L$, and for all $q > 0$ and $t \geq 0$, we have
\begin{equation}\label{eq:Laplace_intro}
\bbE\left[e^{-q\xi_{\rho(t)}}\right]~=~\phi(q)\int_0^{\infty} e^{-\Phi(q)x-\frac{t}{x}} xL(\dd x).
\end{equation}
\end{theo}
\begin{prop}\label{prop:unifBound}
For all $q > 0$, there exists $a(q) \in (0,\infty)$ such that for all $t \geq 0$
\[
\E\left[e^{-q \xi_{\rho(t)}} \right] \leq a(q) \cdot \left(1+t^{1/8}\right) e^{-2 \sqrt{\Phi(q)t}}.
\]
\end{prop}
\begin{rema}
Observe that when $\kappa>0$, $\xi_{\rho(t)}$ may be infinite. Letting $q \downarrow 0$ in~\eqref{eq:Laplace_intro}, we see that
\[
\bbP\big(\xi_{\rho(t)}<\infty\big)=\kappa \int_0^{\infty}e^{-\frac{t}{x}} xL(\dd x), \quad \text{ for all } t \geq 0
\]
(we will see further that the measure $xL(\dd x)$ is finite if and only if $\kappa>0$).
\end{rema}
Under some more precise regularity conditions, we can compute the asymptotic behavior of the Laplace transform of $\xi_{\rho(t)}$ as $t \to \infty$. We first assume that
\begin{equation}\label{hyp:rv}
\tag{$\mathbf{H_{\gamma}}$}
\phi \text{ is regularly varying at 0 with index }\gamma \in (0,1],
\end{equation}
i.e. that the subordinator $\xi$ belongs to the domain of attraction of a stable random variable with index $\gamma$. We recall that a function $f$ is called regularly varying at $0$ with index $\alpha \in \R $ if for all $\lambda > 0$,
\[
\lim_{x \to 0} \frac{f(\lambda x)}{f(x)} = \lambda^\alpha,
\]
and refer to the book of Bingham, Goldie and Teugels~\cite{BGT} for background on regularly varying functions. In particular, observe that under assumption~\eqref{hyp:rv}, we have $\kappa = 0$. In addition to the regular variation of $\phi$, we require an extra technical assumption, which guarantees that the law of $X_1$ is strongly non-lattice, namely that the three following conditions hold
\begin{multline}\label{hyp:tech}
\tag{$\htec$}
\int_0^{\infty}\frac{\pi(\dd x)}{x}=\infty, \\ \int_{1}^{\infty}e^{a\Ree (\Phi(ix))}\dd x<\infty \text{ for some }a>0, \quad \limsup_{x \rightarrow \infty} \Ree (\Phi(ix)) <0.
\end{multline}
These conditions come from our use of results by Doney and Rivero \cite{DR13,DR16_err}, but we are not convinced that they are necessary for the following result to hold.
\begin{theo}\label{thm:asympBehavior}
Assume~\eqref{hyp:rv} and~\eqref{hyp:tech} and define for all $q > 0$
\begin{equation*}\label{eqn:defcq}
b(q) := \frac{\sqrt{\pi}}{(\gamma+1)\Gamma\left(\frac{\gamma}{\gamma + 1}\right)} \cdot \phi(q) \Phi(q)^{\frac{1}{2(\gamma + 1)} - \frac{3}{4}}.
\end{equation*}
Then we have
\[
\E\left[e^{-q \xi_{\rho(t)}}\right]~\underset{t \to \infty} \sim~b(q)\cdot t^{1/4} \Phi^{-1}\left(t^{-1/2}\right) e^{-2 \sqrt{\Phi(q)t}}.
\]
\end{theo}
\section{The time-changed subordinator \texorpdfstring{$\xi_{\rho}$}{xiRho} and its Laplace exponent}\label{sec:Laplace}

In this and the next section, we work with $\xi$ a generic subordinator with Laplace exponent
\[
\phi(q)=\kappa+cq+\int_0^{\infty} \left(1-\exp(-qx)\right)\pi(\dd x), \quad q \geq 0
\]
where $\kappa \geq 0$, $c \geq 0$ and $\pi$ is a measure on $(0,\infty)$ such that $\int_0^{\infty} (1 \wedge x) \pi(\dd x)<\infty$. We assume throughout that $\pi(0,\infty)>0$ and $\xi_0=0$. We recall that $\rho(t)$ is the stopping time defined by
\begin{equation}\label{eq:time_change}
\rho(t)=\inf \left\{u \geq 0 : \int_0^u \xi_r \dd r >t \right\} \text{ for all }t\geq 0.
\end{equation}
Observe that $\rho(0) = \inf\{u \geq 0 : \xi_u > 0\}$ is either null almost surely, if $c > 0$ or $\pi((0,\infty)) = \infty$, or is strictly positive almost surely, when $c=0$ and $\pi$ is finite. In the latter case, it equals the first jump time of $\xi$, therefore $\xi_{\rho(0)}$ is equal to the first jump of $\xi$. Observe also that
\[
\int_0^{\rho(t)} \xi_r \dd r =t \quad \text{for all} \quad t <\int_0^{\bfe(\kappa)}\xi_r \dd r,
\]
\subsection{Definition and first properties of the measure \texorpdfstring{$L$}{L}}\label{sec:L}

We recall from the statement of Theorem~\ref{thm:lapTransform} that $\Phi$ is the primitive of $\phi$ satisfying $\Phi(0) = 0$, i.e.
\begin{align}\label{LK_Phi}
\Phi(q)&=\kappa q+\frac{c}{2}q^2+\int_{\R _+} \left(q-\frac{\left(1-\exp(-qx)\right)}{x}\right)\pi(\dd x) \\
\nonumber &= \kappa q+\frac{c}{2}q^2+ q \pi([1,\infty))+\int_0^{\infty} \left(qx \ind{x\,<\,1} -\left(1-\exp(-qx)\right)\right) \frac{\pi(\dd x)}{x}.
\end{align}
A key remark is that $\Phi$ is the Laplace transform of a Lévy process with diffusion coefficient $c$, drift $\kappa + \pi([1,\infty))$ and a Lévy measure given by the image of $\pi(\dd x)/x$ by $x \mapsto -x$. We denote $X$ such a process, satisfying
\begin{equation*}\label{def_X}
\bbE\left[e^{qX_t}\right]=e^{t \Phi(q)} \quad \text{for all } q,t \geq 0.
\end{equation*}
This class of processes, usually called \emph{spectrally negative Lévy processes}, has remarkable properties and has been studied extensively. We refer to Bertoin's book~\cite[Chapter~VII]{BLevy} for background. Here, we note that either $X$ oscillates ($\limsup_{t\rightarrow \infty} X_t=-\liminf_{t\rightarrow \infty} X_t=+\infty$) when $\kappa=0$, or tends to $+\infty$ ($\lim_{t\rightarrow \infty} X_t=+\infty$) when $\kappa>0$, using that $\Phi'(0)=\phi(0)=\kappa$. It is well known that $\Phi^{-1}:[0,\infty)\mapsto [0,\infty)$ is the Laplace exponent of a subordinator. More precisely,

\begin{lemm}[{\cite[Theorem~1, Chapter~VII]{BLevy}}]
The process $\sigma$ defined for $t \geq 0$ by
\[
\sigma_t = \inf\{ s > 0 : X_s > t\}=\inf\{ u >0:S_u>t\},
\]
where $S_t = \sup_{s\,\leq\,t} X_s$, is a subordinator with Laplace exponent $\Phi^{-1}$. Moreover $S$ is a local time at 0 of the reflected process $S-X$. Therefore, if we let $L$ denote the Lévy measure of $\sigma$, we have
\[
L(\dd x)=n(\zeta \in \dd x)
\]
where $n$ denotes the excursion measure away from 0 of $S-X$ and $\zeta$ the lifetime of an excursion.
\end{lemm}



Since $\Phi^{-1}(0)=0$ and, as $q\rightarrow \infty$, $\Phi^{-1}(q)/q$ converges to $(\kappa+\pi(0,\infty))^{-1}$ when $c=0$ and to $0$ otherwise, the subordinator $\sigma$ has no killing term and a drift equal to $\ind{c=0}(\kappa+\pi(0,\infty))^{-1}$. Consequently the function $\Phi^{-1}$ can be written for all $q\geq 0$ as
\begin{equation}\label{LK_invPhi}
\Phi^{-1}(q)=\frac{q \mathbbm 1_{\{c=0\}}}{\kappa+\pi(0,\infty)}+\int_0^{\infty} \left(1-\exp(-qx)\right)L(\dd x).
\end{equation}
The measure $L$ plays an important role in our study. We underline below a few simple observations. A first one, which will be useful for the proof of Proposition~\ref{prop:unifBound}, is the following connection with a renewal measure.


\begin{lemm}\label{lem:renew}
Let $U$ be the renewal measure of the subordinator $\sigma \circ \xi$ when we consider independent copies of $\xi$ and $\sigma$, and set $V(\dd x):=xL(\dd x)$. Then
\[
U=V + \frac{\mathbbm 1_{\{c=0\}}}{\kappa+\pi(0,\infty)} \delta_0.
\]
Consequently, the function $\bar V$, defined by $\bar V(x):=V([0,x])$ for $x \geq 0$, satisfies:
\[
\bar V(x+y)\leq \bar V(x)+ \bar V(y) + \frac{\mathbbm 1_{\{c=0\}}}{\kappa+\pi(0,\infty)} \qquad \text{for all } x,y \geq 0.
\]
\end{lemm}

\begin{proof}
For background on renewal measures of subordinators we refer to~\cite[Chapter~3]{BLevy} or~\cite[Chapter~5]{KypFluctu}. By differentiating~\eqref{LK_invPhi}, we note that
\[
\frac{1}{\phi \circ \Phi^{-1}(q)}=\left(\Phi^{-1}\right)'(q)=\frac{\mathbbm 1_{\{c=0\}}}{\kappa+\pi(0,\infty)}+ \int_0^{\infty} e^{-qx} V(\dd x).
\]
The function $\phi \circ \Phi^{-1}$ being the Laplace transform of the subordinator $\sigma \circ \xi$ when we consider independent copies of $\xi$ and $\sigma$, this implies that the measure $V$ is the contribution on $(0,\infty)$ of the renewal measure $U$ of $\sigma \circ \xi$ and that $U(\{0\})=\frac{\mathbbm 1_{\{c=0\}}}{\kappa+\pi(0,\infty)}$. It is well-known that the distribution function $x \mapsto U([0,x])$ is then subadditive. This has the particular consequence that the function $\bar V$ satisfies
\[
\bar V(x+y)+ \frac{\mathbbm 1_{\{c=0\}}}{\kappa+\pi(0,\infty)} \leq \bar V(x)+ \frac{\mathbbm 1_{\{c=0\}}}{\kappa+\pi(0,\infty)}+ \bar V(y) + \frac{\mathbbm 1_{\{c=0\}}}{\kappa+\pi(0,\infty)} \quad \text{for all } x,y \geq 0. \qedhere
\]
\end{proof}


We finish with easy integrability properties of $L$. Using that $\Phi^{-1}$ is the Laplace exponent of the subordinator $\sigma$ of first passage times of the Lévy process $X$, we remark that the process $t \mapsto \sigma_t - \frac{t \ind{c = 0}}{\kappa + \pi(0,\infty)}~$ is a compound Poisson process if and only if the diffusion coefficient $c$ of $X$ is null and its Lévy measure is finite, that is
\[
L(0,\infty) < \infty \quad \text{ if and only if } \quad c = 0 \quad \text{ and } \quad \int_0^\infty x^{-1} \pi(\dd  x) <\infty.
\]
Moreover, in this case we have
\[
L(0,\infty) = \lim_{q \to \infty} \left(\Phi^{-1}(q) - \frac{q}{\kappa + \pi(0,\infty)}\right) = \frac{1}{\kappa + \pi(0,\infty)} \int_0^\infty x^{-1} \pi(\dd x).
\]
Similarly, we have $\int_0^{\infty} x L(\dd x) < \infty$ if and only if $\E[\sigma_t] < \infty$ for all $t$. Differentiating~\eqref{LK_invPhi}, we see that
\[
\int_0^{\infty} x L(\dd x) = (\Phi^{-1})'(0) - \frac{\ind{c = 0}}{\kappa + \pi(0,\infty)} = \frac{1}{\kappa} - \frac{\ind{c = 0}}{\kappa + \pi(0,\infty)}.
\]
In particular, $\int_0^{\infty} x L(\dd x) < \infty$ if and only if $\kappa > 0$.
\setcounter{theo}{4}
\subsection{Proof of Theorem~\ref{thm:lapTransform}}\label{sec:proofThmain}

We now prove Theorem~\ref{thm:lapTransform}. In that aim we first introduce, for all $t \geq 0, q>0$,
\begin{equation}\label{def:f}
f(t,q):=(\Phi^{-1})'(q)\bbE\left[e^{-\Phi^{-1}(q) \xi_{\rho(t)}}\right]
\end{equation}
and observe that this bivariate function $f$ is the solution to a simple PDE.

\begin{lemm}\label{lem:easy}
The function $f$ is a strong solution to the equation
\begin{equation}\label{eq:easy}
\partial_q \partial_t f = \partial_t \partial_q f= f.
\end{equation}
Additionally, the function $t \in [0,\infty) \mapsto f(t,q)$ is $\calC^1$ on $(0,\infty)$ and continuous at 0, with $f(0,q)=\int_0^{\infty}xe^{-qx}L(\dd x)$, for all $q>0$.
\end{lemm}

In the proof of this lemma, we will note some regularity properties of $f$ that are a bit more precise than those stated here. Moreover, we will see further with Lemma~\ref{lem:laplaceLaplace} that $f$ is infinitely differentiable on $(0,\infty)^2$.
\setcounter{theo}{7}
Instead of using this approach based on the infinitesimal generator, we prefer to use a more elementary method, which consists in rewriting the function $f$ as follows.
\begin{lemm}\label{lem:rewritten}
For all $t \geq 0$ and $q > 0$, the function $f$ defined in~\eqref{def:f} satisfies
\[
f(t,q) = \int_t^\infty \E\left[\frac{e^{-\Phi^{-1}(q)\xi_{\rho(s)}}}{\xi_{\rho(s)}} \right] \dd s.
\]
\end{lemm}


\begin{proof}
Fix $q>0, t \geq 0$.
We first observe that the change of variable $\rho(s) = u$ leads to
\[
\int_t^\infty \frac{e^{-\Phi^{-1}(q) \xi_{\rho(s)}}}{\xi_{\rho(s)}} \dd s = \int_{\rho(t)}^{\infty} e^{-\Phi^{-1}(q) \xi_u} \dd u.
\]
We underline that when $\kappa > 0$ the subordinator $\xi$ reaches $\infty$ in finite time and that the identity above is valid by using the convention $e^{-\infty} = 0$. Then by Fubini's theorem
\[
\int_t^\infty \E\left[\frac{e^{-\Phi^{-1}(q)\xi_{\rho(s)}}}{\xi_{\rho(s)}} \right] \dd s =
\E\left[\int_{\rho(t)}^{\infty} e^{-\Phi^{-1}(q) \xi_u} \dd u \right].
\]
Applying the strong Markov property for Lévy processes to the stopping time $\rho(t)$, we obtain
\begin{align*}
\E\left[\int_{\rho(t)}^{\infty} e^{-\Phi^{-1}(q) \xi_u} \dd u \right] &= \E\left[e^{-\Phi^{-1}(q) \xi_{\rho(t)}} \int_0^{\infty} e^{-\Phi^{-1}(q) \left(\xi_{u +\rho(t)} - \xi_{\rho(t)}\right)} \dd u \right]\\
&= \E\left[e^{-\Phi^{-1}(q) \xi_{\rho(t)}} \right] \int_0^\infty \E\left[e^{-\Phi^{-1}(q) \xi_{u}} \right] \dd u.
\end{align*}
Observing that $\E[e^{- \Phi^{-1}(q) \xi_t}] = e^{-t \phi(\Phi^{-1}(q))}=e^{-t \Phi'(\Phi^{-1}(q))}$, we conclude that
\[
\int_t^\infty \E\left[\frac{e^{-\Phi^{-1}(q)\xi_{\rho(s)}}}{\xi_{\rho(s)}} \right] \dd s = \frac{1}{\Phi' \circ \Phi^{-1}(q)}\E\left[e^{-\Phi^{-1}(q) \xi_{\rho(t)}} \right] = f(t,q). \qedhere
\]
\end{proof}


\begin{proof}[Proof of Lemma~\ref{lem:easy}]
Using that for all $a,b > 0$, $x \mapsto x^{a} e^{-bx}$ is bounded on $\R_+$, we obtain easily that the function $q \in (0,\infty) \mapsto f(t,q)$ is infinitely differentiable. Moreover, by Lemma~\ref{lem:rewritten}, for all $t \geq 0, q>0$,
\[
\int_t^\infty \E\left[\frac{e^{-\Phi^{-1}(q)\xi_{\rho(s)}}}{\xi_{\rho(s)}} \right] \dd s = f(t,q) < \infty,
\]
therefore $\E[e^{-\Phi^{-1}(q)\xi_{\rho(t)}}/\xi_{\rho(t)}] < \infty$ for almost all $t \geq 0$. Remarking that the function $t \mapsto e^{-\Phi^{-1}(q)\xi_{\rho(t)}}/\xi_{\rho(t)}$ is non-increasing, $t \mapsto \E[e^{-\Phi^{-1}(q)\xi_{\rho(t)}}/{\xi_{\rho(t)}}]$ is a non-increasing function. As a consequence, $t \in (0,\infty) \mapsto f(t,q)$ is a convex decreasing function and differentiable almost everywhere with
\[
\partial_t f(t,q) = -\E\left[\frac{e^{-\Phi^{-1}(q)\xi_{\rho(t)}}}{\xi_{\rho(t)}} \right] \quad \text{for almost all $t > 0$}.
\]
We then observe that for all $t,q > 0$, using Fubini's theorem
\begin{align*}
\int_q^\infty f(t,r) \dd r &= \int_q^\infty (\Phi^{-1})'(r)\E\left[e^{-\Phi^{-1}(r) \xi_{\rho(t)}}\right] \dd r \\
&= \E\left[\int_q^\infty (\Phi^{-1})'(r)e^{-\Phi^{-1}(r) \xi_{\rho(t)}} \dd r \right] = \E\left[\frac{e^{-\Phi^{-1}(q)\xi_{\rho(t)}}}{\xi_{\rho(t)}} \right].
\end{align*}
The function $t \in (0,\infty) \mapsto \E[{e^{-\Phi^{-1}(q)\xi_{\rho(t)}}}/{\xi_{\rho(t)}}]$ is therefore continuous, and consequently the function $t \in (0,\infty) \mapsto f(t,q)$ is $\calC^1$ for all $q > 0$ and
\[
\partial_t f(t,q) = -\int_q^\infty f(t,r) \dd r.
\]
Note that this implies that $q \in (0,\infty) \mapsto \partial_t f(t,q)$ is infinitely differentiable. Then, differentiating the above expression with respect to $q$ shows that $\partial_q \partial_t f(t,q) = f(t,q)$.

Similarly, by Fubini's theorem and the change of variables $\rho(s)=u$,
\begin{align*}
\int_t^\infty f(s,q) \dd s &= (\Phi^{-1})'(q)\bbE\left[\int_t^\infty e^{-\Phi^{-1}(q) \xi_{\rho(s)}} \dd s\right]\\
&= (\Phi^{-1})'(q) \E\left[\int_{\rho(t)}^{\infty} \xi_u e^{-\Phi^{-1}(q) \xi_u} \dd u \right]
\end{align*}
for $t,q>0$. Moreover, we note that for $u> 0$
\begin{multline*}
\E\left[\xi_{\rho(t)+u} e^{-\Phi^{-1}(q) \xi_{\rho(t)+u}}\right]\\
\begin{aligned}
&= \E\left[\left(\xi_{\rho(t)} +\left(\xi_{\rho(t) + u} - \xi_{\rho(t)}\right)\right) e^{-\Phi^{-1}(q) \xi_{\rho(t)}} e^{-\Phi^{-1}(q) \left(\xi_{\rho(t) + u} - \xi_{\rho(t)}\right)} \right]\\
&= \E\left[\xi_{\rho(t)} e^{-\Phi^{-1}(q)\xi_{\rho(t)}} \right] e^{-u \phi\left(\Phi^{-1}(q)\right)} + \E\left[e^{-\Phi^{-1}(q) \xi_{\rho(t)}}\right] u \phi'\left(\Phi^{-1}(q)\right) e^{-u \phi(\Phi^{-1}(q))}.
\end{aligned}
\end{multline*}
This leads to
\begin{multline*}
\int_t^\infty f(s,q) \dd s\\
\begin{aligned}
&= \frac{\left(\Phi^{-1}\right)'(q)}{\phi\left(\Phi^{-1}(q)\right)} \E\left[\xi_{\rho(t)} e^{-\Phi^{-1}(q)\xi_{\rho(t)}} \right] + \phi'\left(\Phi^{-1}(q)\right) \frac{\left(\Phi^{-1}\right)'(q)}{\phi\left(\Phi^{-1}(q)\right)^2} \E\left[e^{-\Phi^{-1}(q) \xi_{\rho(t)}} \right]\\
&= -\partial_q f(t,q),
\end{aligned}
\end{multline*}
since by definition of $f$
\[
\partial_q f(t,q) = (\Phi^{-1})''(q)\E\left[e^{-\Phi^{-1}(q) \xi_{\rho(t)}}\right] - (\Phi^{-1})'(q)^2 \E\left[\xi_{\rho(t)} e^{-\Phi^{-1}(q)\xi_{\rho(t)}} \right]
\]
and $(\Phi^{-1})' = 1/\phi (\Phi^{-1})$. Differentiating with respect to $t$ shows that we also have $\partial_t \partial_q f(t,q) = f(t,q)$.

The continuity at 0 of $t \mapsto f(t,q)$ is obvious, for all $q>0$, by using the expression of $f$ as an integral as noticed in Lemma~\ref{lem:rewritten}. It remains to identify the initial expression for $f(0,q)$ in terms of the measure $L$. Since $\xi_{\rho(0)}=0$ when $c>0$ or $\pi$ is infinite, and $\xi_{\rho(0)}$ is the first jump time of $\xi$ otherwise, one has
\[
\bbE\left[e^{-q \xi_{\rho(0)}}\right]=\left\{
\begin{array}{ll} 1 & \text{when $c>0$ or $\pi$ is infinite}\\ \frac{1}{\kappa+\pi(0,\infty)}\int_0^{\infty}e^{-qx} \pi(\dd x) & \text{when $c=0$ and $\pi$ is finite.}
\end{array}\right.
\]
When $c>0$ or $\pi$ is infinite, recalling~\eqref{LK_invPhi}, this indeed yields
\[
f(0,q)=(\Phi^{-1})'(q)=\int_0^{\infty} xe^{-qx} L(\dd x).
\]
When $c=0$ and $\pi$ is finite, we have
\[
f(0,q)=(\Phi^{-1})'(q) \frac{1}{\kappa+\pi(0,\infty)}\int_0^{\infty}e^{-\Phi^{-1}(q)x} \pi(\dd x).
\]
Using that in this case $\Phi(q)=(\kappa+\pi(0,\infty))q-\int_0^{\infty}(1-e^{-qx})\pi(\dd x)/x$, we see that
\[
q=(\kappa+\pi(0,\infty))\Phi^{-1}(q)-\int_0^{\infty}\left(1-e^{-\Phi^{-1}(q)x}\right)\pi(\dd x)/x,
\]
which, by differentiating, leads us to
\[
f(0,q)=(\Phi^{-1})'(q) \frac{1}{\kappa+\pi(0,\infty)}\int_0^{\infty}e^{-\Phi^{-1}(q)x} \pi(\dd x)=(\Phi^{-1})'(q) -\frac{1}{\kappa+\pi(0,\infty)}.
\]
By~\eqref{LK_invPhi} and since $c=0$, this last expression is equal to $\int_0^{\infty} xe^{-qx} L(\dd x)$ as expected.
\end{proof}


Solving the partial differential equation of Lemma~\ref{lem:easy}, we obtain an expression for the Laplace transform of the function $t \mapsto f(t,q)$ for each $q > 0$. This leads to an explicit expression of $f(t,q)$ in terms of the measure $L$, which proves Theorem~\ref{thm:lapTransform}.

\begin{lemm}\label{lem:laplaceLaplace}
For $\lambda>0$, we write $F_{\lambda} : q \in (0,\infty)\mapsto \int_0^{\infty} e^{-\lambda t} f(t,q) \dd t$. Then for all fixed $q>0$, we have
\begin{equation}\label{eq:Laplace}
F_{\lambda}(q)=\int_0^{\infty} \left(\int_0^{\infty} e^{-qx -\frac{t}{x}} xL(\dd x)\right) e^{-\lambda t} \dd t, \quad \forall \lambda>0
\end{equation}
which implies that $f(t,q)= \int_0^{\infty} e^{-qx -\frac{t}{x}} xL(\dd x)$ for all $t \geq 0, q>0$. In particular, $f$ is infinitely differentiable on $(0,\infty)^2$.
\end{lemm}


\begin{proof}
Fix $\lambda > 0$ and note that $F_{\lambda}(q)$ is finite for all $q > 0$, since the function $t \in (0,\infty) \mapsto f(t,q)$ is bounded. An integration by part gives
\begin{align*}
F_{\lambda}(q)=& \left[-\frac{e^{-\lambda t}}{\lambda}f(t,q) \right]_{t=0}^{t=\infty} + \frac{1}{\lambda} \int_0^{\infty}e^{-\lambda t} \partial _t f(t,q) \dd t \\
=& \frac{1}{\lambda} \int_0^{\infty} xe^{-qx} L(\dd x)+\frac{1}{\lambda} \int_0^{\infty}e^{-\lambda t} \partial _t f(t,q) \dd t,
\end{align*}
using that $e^{-\lambda t}f(t,q) \rightarrow 0 $ as $t\rightarrow \infty$, and that $f(t,q) \rightarrow \int_0^{\infty} xe^{-qx} L(\dd x)$ as $t \rightarrow 0$ by Lemma~\ref{lem:easy}. We then have, for all $q>0$,
\begin{align*}
\partial_q F_{\lambda}(q)&= \frac{-1}{\lambda} \int_0^{\infty} x^2e^{-qx} L(\dd x)+ \frac{1}{\lambda} \int_0^{\infty}e^{-\lambda t} \partial _q\partial _t f(t,q) \dd t \\
&= \frac{-1}{\lambda} \int_0^{\infty} x^2e^{-qx} L(\dd x)+ \frac{1}{\lambda} F_{\lambda}(q),
\end{align*}
using again Lemma~\ref{lem:easy} (we can apply Lebesgue's dominated convergence theorem since, for all $q_0>0$ and $t>0$,
\[
\sup_{q\,\geq\,q_0} \left|e^{-\lambda t}\partial_q \partial_t f(t,q)\right|=\sup_{q\,\geq\,q_0} \left|e^{-\lambda t} f(t,q)\right|=e^{-\lambda t} f(t,q_0),
\]
since $(\Phi^{-1})'$ is decreasing and $\Phi^{-1}$ is increasing).

The function $q \in (0,\infty) \mapsto F_{\lambda}(q)$ is therefore a solution to a linear differential equation of the first order. By standard arguments, it follows that
\[
F_{\lambda}(q)=e^{\frac{q}{\lambda}}\int_q^{\infty} e^{-\frac{u}{\lambda}} \left(\frac{1}{\lambda} \int_0^{\infty} x^2e^{-ux} L(\dd x) \right) \dd u.
\]
Using Fubini's theorem and the change of variable $t=(u-q)x/\lambda$, we then get the expression~\eqref{eq:Laplace} for all $q,\lambda>0$. The final equation follows from Laplace inversion theorem and the fact that the functions $t\mapsto f(t,q)$ and $t \mapsto \int_0^{\infty} e^{-qx -\frac{t}{x}} xL(\dd x)$ are continuous on $(0,\infty)$.
\end{proof}


\begin{proof}[{Proof of Theorem~\ref{thm:lapTransform}}]
Observe that the definition~\eqref{def:f} of $f$ implies that
\[
f(\Phi(q),t) = (\Phi^{-1})'(\Phi(q)) \E\left[e^{-q \xi_{\rho(t)}} \right] = \frac{1}{\phi(q)}\E\left[e^{-q \xi_{\rho(t)}} \right].
\]
As a result, Lemma~\ref{lem:laplaceLaplace} yields
\[
\bbE\left[e^{-q\xi_{\rho(t)}}\right] = \phi(q) f\left(t,\Phi(q)\right)=\phi(q) \int_0^{\infty} e^{-\Phi(q)x -\frac{t}{x}} \ xL(\dd x), \quad \forall t\geq 0, q>0,
\]
which completes the proof.
\end{proof}
\subsection{Proof of Proposition~\ref{prop:unifBound}}\label{sec:bounds}

The proof of Proposition~\ref{prop:unifBound} relies on Theorem~\ref{thm:lapTransform} and Lemma~\ref{lem:renew}. In fact, we use the forthcoming Lemma~\ref{eq:dev}, an immediate corollary of Theorem~\ref{thm:lapTransform}, which rewrites in a more convenient way the expression of the Laplace transform of $\xi_{\rho(t)}$. From this lemma and Lemma~\ref{lem:renew}, we note that for all $t \geq 0,q>0$
\begin{multline}\label{eq:largebound}
\bbE\left[e^{-q\xi_{\rho(t)}}\right] \\
\leq \phi(q)
e^{-2 \sqrt{\Phi(q)t}}
\int_0^\infty 2u e^{-u^2}
\left(\bar V\left(2
\left(\frac{\sqrt{u^4+4u^2 \sqrt{\Phi(q) t}}}
{2 \Phi(q)} \right)\right)
+ \frac{\mathbbm 1_{\{c=0\}}}
{\kappa+\pi(0,\infty
)} \right)\dd u
,
\end{multline}
where we recall that $\bar V(x)=\int_{[0,x]}u L(\dd u)$. It remains to bound $\bar V$ from above to conclude. Since $\phi$ is concave, the function $q \mapsto \phi(q)/q$ is decreasing on $(0,\infty)$ and $\lim_{q \downarrow 0} \phi(q)/q \in (0,\infty]$. Consequently there exists $c_1 \in (0,\infty)$ such that for all $q \in [0,1]$, we have $\Phi^{-1}(q) \leq c_1 \sqrt q$. Then write for $x \geq 1$,
\[
\bar V(x) = x\int_{[0,x]} \frac{u}{x} L(\dd u) \leq 3 x \int_{[0,x]} \left(1-e^{- \frac{u}{x}} \right) L(\dd u) \leq 3 x \Phi^{-1}(1/x) \leq 3c_1 \sqrt x
\]
where we have used for the first inequality that $y \leq 3(1-e^{-y})$ for $y \in [0,1]$ and~\eqref{LK_invPhi} for the second inequality. This leads to the existence of $c_2 \in (0,\infty)$ such that
\[
\bar V(x) \leq c_2 \left(1 + \sqrt x\right), \quad \forall x \geq 0,
\]
which in turn implies that for all $u,t \geq 0$, $q>0$,
\[
\bar V\left(2 \left(\frac{\sqrt{u^4+4u^2 \sqrt{\Phi(q) t}}}{2 \Phi(q)} \right)\right) \leq c_2 \left(1 + \frac{u+\sqrt{2u} \ \Phi(q)^{1/8}t^{1/8}}{\sqrt{\Phi(q)}}\right).
\]
And then, with~\eqref{eq:largebound}, we get the upper bound
\[
\bbE\left[e^{-q\xi_{\rho(t)}}\right] \leq c_3\left(1+(\Phi(q))^{-1/2}+(\Phi(q))^{-3/8}t^{1/8} \right) \phi(q)e^{-2 \sqrt{\Phi(q)t}}
\]
for some $c_3>0$, as expected.
\section{Asymptotics of \texorpdfstring{$\xi_{\rho}$}{xiRho}}\label{sec:asymptotic_sub}

Theorem~\ref{thm:lapTransform} will allow us to obtain the asymptotic behavior of $\bbE[\exp(-q\xi_{\rho(t)})]$ as $t\rightarrow \infty$, both when $q$ is fixed or when $q=q(t)$ is an appropriate function of $t$ chosen to obtain the scaling limit of $\xi_{\rho(t)}$. Indeed, observe that in the expression
\begin{equation}\label{eq:Laplacebis}
\E\left[e^{-q \xi_{\rho(t)}}\right] = \phi(q)\int_0^\infty e^{-\Phi(q)x-\frac{t}{x}} xL(\dd x)
\end{equation}
the function $x \mapsto e^{-\Phi(q)x-\frac{t}{x}}$ attains its maximum at $x=\sqrt{t/\Phi(q)}$, which is equal to
\[
e^{-2 \sqrt{\Phi(q)t}}.
\]
Therefore, as long as $L$ is regular enough that we may apply Laplace's saddle point method, we expect the leading term in the asymptotic behavior of $\E[e^{-q \xi_{\rho(t)}}]$ to be $e^{-2 \sqrt{\Phi(q)t}}$, with a polynomial correction. The first aim of this section is to prove Theorem~\ref{thm:asympBehavior}, i.e. that under conditions~\eqref{hyp:rv} and~\eqref{hyp:tech}, an equivalent of $\E[e^{-q \xi_{\rho(t)}}]$ as $t \to \infty$ can be computed explicitly. Then, with very similar computations, we will be able to obtain the scaling limit of $\xi_{\rho(t)}$ as $t \to \infty$ and prove Theorem~\ref{thm:asymp_distribution}.

The starting point of the proofs of Theorem~\ref{thm:asympBehavior} and Theorem~\ref{thm:asymp_distribution} rely on the following alternative expression for the Laplace transform~\eqref{eq:Laplacebis}. Let us recall the notation $V(\dd x) = x L(\dd x)$, and for $x \geq 0$
\[
\bar{V}(x) = V([0,x]) \quad \text{and} \quad \bar{L}(x) :=L(x,{\infty}).
\]


\begin{lemm}\label{eq:dev}
For $u, t, q >0$, let
\[
a_{u,t,q} = \frac{u^2+2 \sqrt{\Phi(q)t}}{2\Phi(q)} \quad \text{and} \quad b_{u,t,q} = \frac{\sqrt{u^4+4u^2 \sqrt{\Phi(q)t}}}{2\Phi(q)},
\]
\emph{(}note that $a_{u,t,q}>b_{u,t,q}$\emph{)}\footnote{More precisely, we have $(2\Phi(q) a_{u,t,q})^2 - (2 \Phi(q)b_{u,t,q})^2 = 4 \Phi(q) t$. In particular, $2 \Phi(q)(a_{u,t,q} - b_{u,t,q}) \to \infty$ as long as $4 \Phi(q)t \to \infty$, uniformly in $u$.}. We then have
\begin{equation*}
\E\left[e^{-q\xi_{\rho(t)}}\right] = \phi(q) e^{-2 \sqrt{\Phi(q)t}}\int_0^\infty 2u e^{-u^2} \left(\bar V \left(a_{u,t,q}+b_{u,t,q}\right) - \bar V \left(a_{u,t,q}-b_{u,t,q}\right) \right) \dd u.
\end{equation*}
\end{lemm}


\begin{proof}
Since $\Phi(q)x + \frac{t}{x} = 2 \sqrt{\Phi(q)t} + (\sqrt{\Phi(q)x}-\sqrt{\frac{t}{x}})^2$, we can rewrite~\eqref{eq:Laplacebis} as
\begin{align*}
\frac{\E\left[e^{-q\xi_{\rho(t)}}\right]}{\phi(q) e^{-2 \sqrt{\Phi(q)t}}} &= \int_0^\infty e^{-\left(\sqrt{\Phi(q)x}-\sqrt{\frac{t}{x}}\right)^2} x L(\dd x)\\
&= \int_0^{\infty} \int_{\left|\sqrt{\Phi(q)x}-\sqrt{\frac{t}{x}}\right|}^{\infty}2ue^{-u^2} \dd u \ xL(\dd x)\\
&= \int_0^{\infty} 2ue^{-u^2} \Bigg(\int_0^{\infty} \mathbbm 1_{\Big\{\Big|\sqrt{\Phi(q)x}-\sqrt{\frac{t}{x}}\Big| \leq u\Big\}} xL(\dd x)\Bigg) \dd u,
\end{align*}
by Fubini's theorem. Note then that
\[
\left\{x \in \R : \left|\sqrt{\Phi(q)x}-\sqrt{\frac{t}{x}}\right| \leq u\right\} = \big[a_{u,t,q}-b_{u,t,q},a_{u,t,q}+b_{u,t,q}\big]
\]
and that the Lebesgue measure of $\{u>0:V(\{a_{u,t,q}-b_{u,t,q}\})>0\}$ is null to complete the proof of Lemma~\ref{eq:dev}.
\end{proof}



The function $\bar V$ is regularly varying at $\infty$ under~\eqref{hyp:rv} which is important but not sufficient for our purpose, as we have to control the increments of $\bar V$. To that end, we apply a result of Doney and Rivero~\cite{DR13} in the next section, which relies on~\eqref{hyp:tech}. We then complete the proofs of Theorem~\ref{thm:asympBehavior} in Section~\ref{sec:prooffix} and Theorem~\ref{thm:asymp_distribution} in Section~\ref{sec:proofscaling}.


\subsection{Preliminaries: regular variation and asymptotic behavior of the increments of \texorpdfstring{$\bar V$}{V}}\label{sec:incrementsofV}

We first study the regular variation of $\bar V$ and $\bar{L}$.


\begin{lemm}\label{lem:varVL}
If $\phi$ is regularly varying at $0$ with index $\gamma \in [0,1]$, the function $\bar V$ is regularly varying at $\infty$ with index $\frac{\gamma}{\gamma+1} \in [0,1/2]$. More precisely:
\begin{enumerate}\romanenumi
\item\label{lemm3.2.1} When $\gamma>0$, $\bar L$ is also regularly varying at $\infty$, with index $\frac{-1}{\gamma+1} \in (-1,-1/2]$, and
\[
\bar L(x) \underset{x \rightarrow \infty}{\sim} \frac{1}{\Gamma\left(\frac{\gamma}{1+\gamma}\right)} \cdot \Phi^{-1}\left(\frac{1}{x}\right) \underset{x \rightarrow \infty}{\sim} \frac{\gamma \bar V(x)}{x}.
\]
\item\label{lemm3.2.2} When $\gamma=0$, $\bar L(x) \ll \bar V(x)/x$ and we still have the above behavior for $\bar V$ when moreover $\kappa=0$:
\[
\bar V(x) \underset{x \rightarrow \infty}{\sim} x \Phi^{-1}\left(\frac{1}{x}\right)
\]
whereas when $\kappa>0$,
\[
\bar V(x) \underset{x \rightarrow \infty}{\rightarrow} \frac{1}{\kappa}- \frac{\mathbbm 1_{\{c=0\}}}{\kappa +\pi(0,\infty)}.
\]
\end{enumerate}
\end{lemm}

These results on the links between the asymptotics of the tail of the Lévy measure of a subordinator and its Laplace exponent are classical. We prove them quickly for the sake of completeness and to integrate the function $\bar V$. Note that $x \Phi^{-1}\left(\frac{1}{x}\right) \rightarrow \frac{1}{\kappa}$ as $x \rightarrow \infty$ when $\kappa>0$, so in general the two assertions of~\eqref{lemm3.2.2} cannot be merged.


\begin{proof}
We rewrite~\eqref{LK_invPhi} as
\begin{equation}\label{eq:rephrased}
\begin{split}
\Phi^{-1}(q) &=\frac{q \mathbbm 1_{\{c=0\}}}{\kappa+\pi(0,\infty)}+\int_0^q \left(\int_0^{\infty} e^{-xu}V(\dd x) \right) \dd u\\
&=\frac{q \mathbbm 1_{\{c=0\}}}{\kappa+\pi(0,\infty)}+q\int_0^{\infty}e^{-qu} \bar L(u) \dd u,
\end{split}
\end{equation}
using Fubini's theorem.

We first assume that $\kappa = 0$. In this case, $\Phi(q)/q \to 0$ as $q \to 0$, therefore $\Phi^{-1}(q) \gg q$ when $q \rightarrow 0$. Consequently,
\begin{equation}\label{eq:miseenplace}
\Phi^{-1}(q) \underset{q \rightarrow 0}{\sim} \int_0^q \left(\int_0^{\infty} e^{-xu}V(\dd x) \right) \dd u.
\end{equation}
As $\phi$ is regularly varying at $0$ with index $\gamma$, $\Phi^{-1}$ is regularly varying at $0$ with index $(\gamma+1)^{-1}$, therefore so is the integral on the right-hand side. Since the function $u \mapsto \int_0^{\infty} e^{-xu}V(\dd x)$ is monotone (decreasing), it is also regularly varying at $0$, with index $(\gamma+1)^{-1}-1$, by the monotone density theorem (\cite[Theorem~1.7.2b]{BGT}) and
\[
\int_0^{\infty} e^{-xq}V(\dd x) \underset{q \rightarrow 0}{\sim} \frac{q^{-1}}{\gamma+1} \int_0^q \left(\int_0^{\infty} e^{-xu}V(\dd x) \right) \dd u \underset{q \rightarrow 0}{\sim} \frac{q^{-1}}{\gamma+1} \Phi^{-1}(q).
\]
It is then sufficient to use Karamata's Tauberian theorem (\cite[Theorem~1.7.1]{BGT}) to deduce from this the regular variation of $\bar V$ and the equivalence settled in the statements~\eqref{lemm3.2.1} and~\eqref{lemm3.2.2} of the lemma, regarding $\bar V$.

The behavior of $\bar{L}$ is obtained in a similar fashion. Using~\eqref{eq:rephrased} and that $\Phi^{-1}(q) \gg q$ as $q \to 0$, we have
\[
\Phi^{-1}(q)\underset{q \rightarrow 0}{\sim} q \int_0^{\infty}e^{-qu} \bar{L}(u) \dd u.
\]
Hence, applying Karamata's tauberian theorem we obtain that
\[
\int_0^{x} \bar L(u) \dd u \underset{x \rightarrow \infty}{\sim} \frac{1}{\Gamma\left(1+\frac{\gamma}{1+\gamma}\right)} \cdot x \Phi^{-1}(1/x),
\]
those two functions being regularly varying with index $\gamma/(1+\gamma)$. And then applying the monotone density theorem to get the remaining parts of assertions~\eqref{lemm3.2.1} and~\eqref{lemm3.2.2}, still when $\kappa=0$.

We now turn to the case $\kappa > 0$. In this situation
\[
\lim_{x \to \infty} \bar{V}(x) = \int_0^\infty y L(y) \dd y =\frac{1}{\kappa}- \frac{\mathbbm 1_{\{c=0\}}}{\kappa +\pi(0,\infty)},
\]
as noticed at the end of Section~\ref{sec:L}.
\end{proof}


To get some information on the asymptotic behavior of the increments of $\bar V$, or equivalently of $\bar L$, we recall from Section~\ref{sec:Laplace} that the measure $L$ can be interpreted as the image measure of the excursion measure, noted $n$, of a spectrally negative Lévy process $X$ with Laplace exponent $\Phi$ reflected below its maximum by the application $\zeta$ associating to an excursion its lifetime. In other words, we have $\bar L(x)=n(\zeta>x)$. We can therefore apply a result of Doney and Rivero~\cite{DR13,DR16_err} to obtain the following local properties.

\begin{lemm}\label{lem:increments}
Assume~\eqref{hyp:rv} and~\eqref{hyp:tech}.
\begin{enumerate}\romanenumi
\item\label{lemm3.3.1} For any $\Delta(x)>0$ such that $\Delta(x)=o(x)$ as $x \rightarrow \infty$,
\[
\sup_{\Delta\,\in\,(0,\Delta(x)]} \left|\frac{x}{\bar L(x)} \cdot \frac{\bar L(x)-\bar L(x+\Delta)}{\Delta} - \frac{1}{1+\gamma}\right| \underset{x \rightarrow \infty}\rightarrow 0.
\]
\item\label{lemm3.3.2} There exists $x_0>0$ such that
\[
\sup_{\Delta\,>\,0} \left|\bar L(x)-\bar L(x+\Delta) \right| \leq \frac{2 \Delta \bar L(x)}{x} \quad \text{for all } x \geq x_0.
\]
\end{enumerate}
\end{lemm}

\begin{proof}
By~\cite[Theorem~1$\MK$(iii)]{DR13} and the erratum~\cite{DR16_err}, under the hypotheses~\eqref{hyp:rv} and~\eqref{hyp:tech},
\[
\sup_{\Delta\,\in\,(0,\Delta_0]} \left|\frac{x}{\bar L(x)} \cdot \frac{\bar L(x)-\bar L(x+\Delta)}{\Delta} - \frac{1}{1+\gamma}\right| \underset{x \rightarrow \infty}\rightarrow 0
\]
for any fixed $\Delta_0>0$ (note that with the notation $\alpha,\rho$ of~\cite{DR13}, we have here $\alpha=1+\gamma$ and $\bar \rho=\frac{1}{1+\gamma}$). Now fix $\varepsilon>0$ and let $x_{\varepsilon}>0$ be such that for $x \geq x_{\varepsilon}$
\[
\sup_{\Delta\,\in\,(0,1]} \left|\frac{x}{\bar L(x)} \cdot \frac{\bar L(x)-\bar L(x+\Delta)}{\Delta} - \frac{1}{1+\gamma}\right| \leq \varepsilon.
\]
Then for any $\Delta>0$ and all $x \geq x_{\varepsilon}$
\begin{align*}
\bar L(x)-\bar L(x+\Delta) &= \sum_{k=0}^{\lfloor\Delta \rfloor-1} \left(\bar L(x+k)-\bar L(x+k+1)\right)+\bar L(x+\lfloor\Delta \rfloor)-\bar L(x+\Delta) \\
& \leq \left(\frac{1}{1+\gamma} + \varepsilon\right) \left(\sum_{k=0}^{\lfloor\Delta \rfloor-1} \frac{\bar L(x+k)}{x+k} + \frac{\bar L(x+\lfloor\Delta \rfloor)}{x+\lfloor\Delta \rfloor} \cdot (\Delta-\lfloor\Delta \rfloor)\right) \\
& \leq \left(\frac{1}{1+\gamma} + \varepsilon\right) \frac{\bar L(x)}{x} \cdot \Delta
\end{align*}
where for the last inequality we use that $x \mapsto \bar L(x)/x$ is decreasing. This implies the point~\eqref{lemm3.3.2} of the lemma and is a first step in the proof of the point~\eqref{lemm3.3.1}. To complete the proof of the uniform convergence on intervals of the form $(0,\Delta(x)]$ with $\Delta(x)=o(x)$, we note, similarly as above, that for any $\Delta \in (0,\Delta(x)]$ and all $x \geq x_{\varepsilon}$,
\begin{align*}
\bar L(x)-\bar L(x+\Delta) &\geq \left(\frac{1}{1+\gamma} - \varepsilon\right) \left(\sum_{k=0}^{\lfloor\Delta \rfloor-1} \frac{\bar L(x+k)}{x+k} + \frac{\bar L(x+\lfloor\Delta \rfloor)}{x+\lfloor\Delta \rfloor} \cdot (\Delta-\lfloor\Delta \rfloor)\right) \\
&\geq
\left(\frac{1}{1+\gamma} - \varepsilon\right) \frac{\bar L(x+\Delta(x))}{x+\Delta(x)} \cdot \Delta.
\end{align*}
Now, since the function $x \mapsto \bar L(x)/x$ is regularly varying as $x \rightarrow \infty$ with index $-\frac{1}{1+\gamma}-1$, by the Uniform Convergence Theorem (\cite[Theorem~1.5.2]{BGT}),\linebreak for all $x$ large enough
\[
\frac{\bar L(x+\Delta(x))}{x+\Delta(x)} \cdot \frac{x}{\bar L(x)} \geq \frac{1}{\left(1+\frac{\Delta(x)}{x}\right)^{\frac{1}{1+\gamma}+1}}-\varepsilon.
\]
Taking $x$ larger if necessary, we may assume that this last lower bound is in turn larger than $1-2 \varepsilon$ since $\Delta(x)=o(x)$. So, finally we have that for all $x$ large enough and then all $\Delta \in (0,\Delta(x)]$,
\[
\bar L(x)-\bar L(x+\Delta) \geq \left(\frac{1}{1+\gamma} - \varepsilon\right) (1-2\varepsilon) \frac{\bar L(x)}{x} \cdot \Delta
\]
as expected.
\end{proof}


This readily yields the following uniform behavior of the increments of $V$.


\begin{coro}\label{cor:increments}
Assume~\eqref{hyp:rv} and~\eqref{hyp:tech}.
\begin{enumerate}\romanenumi
\item\label{coro3.4.1} For any $\Delta(x)>0$ such that $\Delta(x)=o(x)$ as $x \rightarrow \infty$,
\[
\sup_{\Delta\,\in\,(0,\Delta(x)]} \left|\frac{\bar V(x+\Delta)-\bar V(x)}{\Delta \bar L(x)} - \frac{1}{1+\gamma}\right| \underset{x \rightarrow \infty}\rightarrow 0.
\]
\item\label{coro3.4.2} There exists $x_1>0$ such that
\[
\sup_{\Delta\,>\,0} \left|\bar V(x)-\bar V(x+\Delta) \right| \leq 2 \Delta \bar L(x) +2 \frac{\bar L(x)}{x} \Delta^2 \quad \text{for all } x \geq x_1.
\]
\end{enumerate}
\end{coro}

\begin{proof}
Recall that $V(\dd x)=x L(\dd x)$, so that
\begin{equation}\label{eq:VL}
\bar V(x+\Delta)-\bar V(x)=x\left(\bar L(x)-\bar L(x+\Delta)\right)+\int_0^{\Delta} \left(\bar L(x+y)- \bar L(x+\Delta)\right) \dd y.
\end{equation}
From Lemma~\ref{lem:increments}$\MK$\eqref{lemm3.3.1}, we have on the one hand that
\[
\sup_{\Delta\,\in\,(0,\Delta(x)]} \left|\frac{x\left(\bar L(x)-\bar L(x+\Delta)\right)}{\Delta \bar L(x)}- \frac{1}{1+\gamma}\right| \underset{x \rightarrow \infty}\rightarrow 0
\]
and on the other hand that for $x$ large enough and then all $\Delta \in (0,\Delta(x)]$
\begin{align*}
\left|\int_0^{\Delta} \left(\bar L(x+y)- \bar L(x+\Delta)\right) \dd y \right| &\leq \frac{2}{1+\gamma} \int_0^{\Delta} \frac{\bar L(x+y)}{x+y} (\Delta-y) \dd y \\
&\leq \frac{1}{1+\gamma} \cdot \frac{\bar L(x)}{x} \cdot \left(\Delta\right)^2,
\end{align*}
hence
\[
\sup_{\Delta\,\in\,(0,\Delta(x)]} \frac{\left|\int_0^{\Delta} \left(\bar L(x+y)- \bar L(x+\Delta)\right) \dd y \right|}{\Delta \bar L(x)} \leq \frac{1}{1+\gamma} \cdot \frac{\Delta(x)}{x} \longrightarrow 0.
\]
This gives~\eqref{coro3.4.1}. The upper bound in~\eqref{coro3.4.2} is obtained in a similar way, by combining Lemma~\ref{lem:increments}$\MK$\eqref{lemm3.3.2} with~\eqref{eq:VL}.
\end{proof}

\subsection{Asymptotics of \texorpdfstring{$\bbE[\exp(-q\xi_{\rho(t)})]$ when $t \rightarrow \infty$}{the Laplace transform at large times}}\label{sec:prooffix}

In this section $q>0$ is fixed and we let $t \rightarrow \infty$. We have set up all the key elements necessary for the proof of Theorem~\ref{thm:asympBehavior}, which ends as follows.
\begin{proof}[Proof of Theorem~\ref{thm:asympBehavior}.] Set for $u,t,q>0$,
\[
I_{u,t,q}:=\bar V \left(a_{u,t,q}+b_{u,t,q}\right)-\bar V \left(a_{u,t,q}-b_{u,t,q}\right),
\]
with the notation $a,b$ of Lemma~\ref{eq:dev}, so that
\[
\bbE\left[e^{-q\xi_{\rho(t)}}\right]=\phi(q)e^{-2 \sqrt{\Phi(q)t}}\int_0^{\infty} 2ue^{-u^2} I_{u,t,q} \dd u.
\]
Then first note that for each fixed $u,q>0$,
\[
a_{u,t,q} \underset{t \rightarrow \infty}{\sim} \frac{\sqrt t}{\sqrt{\Phi(q)}} \quad \text{and} \quad b_{u,t,q} \underset{t \rightarrow \infty}{\sim} \frac{u t^{1/4}}{(\Phi(q))^{3/4}}
\]
and that Corollary~\ref{cor:increments}$\MK$\eqref{coro3.4.1} yields
\[
I_{u,t,q} \underset{t \rightarrow \infty}{\sim} \frac{2 b_{u,t,q} \bar L(a_{u,t,q}-b_{u,t,q})}{1+\gamma}.
\]
Together with the regular variation of $\bar L$ at $\infty$ (Lemma~\ref{lem:varVL}$\MK$\eqref{lemm3.2.1}), we get that
\[
\frac{I_{u,t,q}}{t^{1/4} \bar L(\sqrt t)} \underset{t\rightarrow \infty} \longrightarrow \frac{2u}{1+\gamma} \cdot \Phi(q)^{\frac{1}{2(1+\gamma)}-\frac{3}{4}}.
\]
It remains to conclude using the dominated converge theorem that this implies
\begin{equation}\label{cv:DCT}
\frac{\int_0^{\infty}2ue^{-u^2} I_{u,t,q}\dd u}{t^{1/4} \bar L(\sqrt t)} \underset{t\rightarrow \infty} \longrightarrow \frac{1}{1+\gamma} \cdot \Phi(q)^{\frac{1}{2(1+\gamma)}-\frac{3}{4}} \int_0^{\infty}4u^2e^{-u^2}\dd u.
\end{equation}
The convergence indeed gives Theorem~\ref{thm:asympBehavior} since $\Gamma(\frac{\gamma}{\gamma+1})\bar L(\sqrt t) \sim_{t \rightarrow \infty} \Phi^{-1}(1/\sqrt t)$ by Lemma~\ref{lem:varVL}$\MK$\eqref{lemm3.2.1} and $\int_0^{\infty} 4u^2e^{-u^2}\dd u=\sqrt \pi$.


To apply the dominated convergence theorem and get~\eqref{cv:DCT}, we need to split the term under the integral into two pieces, according to whether $u\leq u_t:=\left(\Phi(q)t \right)^{1/4}/2~$ or $u>u_t$. First, for $u \leq u_t$, we note that
\[
a_{u,t,q}-b_{u,t,q} \geq \frac{2 \sqrt{\Phi(q)t}-2u \left(\Phi(q)t\right)^{1/4}}{2 \Phi(q)} \geq \frac{\sqrt t}{2\sqrt{\Phi(q)}}.
\]
So, by Corollary~\ref{cor:increments}$\MK$\eqref{coro3.4.2}, we have that for $t$ large enough and all $u\leq u_t$,
\begin{align*}
\left|I_{u,t,q}\right| \leq& 4 b_{u,t,q} \bar L\left(a_{u,t,q}-b_{u,t,q}\right) \left(1+ \frac{2b_{u,t,q}}{a_{u,t,q}-b_{u,t,q}} \right) \\
\leq& 4 b_{u,t,q} \left(1+ \frac{2\Phi(q)b_{u,t,q}(a_{u,t,q}+b_{u,t,q})}{t} \right) \bar L \left({\frac{\sqrt t}{2\sqrt{\Phi(q)}}}\right)
\end{align*}
using that $\bar{L}$ is non-increasing and that $(a_{u,t,q}-b_{u,t,q})(a_{u,t,q}+b_{u,t,q})=t/\Phi(q)$. Besides, as a consequence of the regular variation of $\bar L$ with index $-1/(1+\gamma)$, we have that for all $x \geq y$ both sufficiently large,
\[
\frac{\bar L(y)}{\bar L(x)} \leq 2 \left(\frac{x}{y}\right)^{\frac{2}{(1+\gamma)}},
\]
applying classical Potter's bounds for regularly varying functions (\cite[Theorem~1.5.6]{BGT}). Therefore, taking $t$ larger if necessary, we also have that
\[
\bar L\left({\frac{\sqrt t}{2\sqrt{\Phi(q)}}}\right)\leq c_1(q)\bar L(\sqrt t))
\]
for a constant $c_1(q)$ depending only on $q$ (using Potter's bounds when
$
2 \sqrt{\Phi(q)}\geq 1
$
and that $\bar L$ is non-increasing when $2\sqrt{\Phi(q)}\leq 1$).
Using the definitions of $a_{u,t,q},b_{u,t,q}$, we then deduce from these bounds that for $t$ large enough and all $u \leq u_t$,
\begin{equation*}\label{majo:I}
\left|I_{u,t,q}\right|\leq c_2(q) t^{1/4} \bar L\big({\sqrt t}\big) P(u),
\end{equation*}
where $P$ is a polynomial and $c_2(q)$ does not depend on $t$ and $u\leq u_t$. We can therefore apply the dominated convergence theorem and get that
\[
\frac{\int_0^{\infty} 2ue^{-u^2} I_{u,t,q} \ind{u\,\leq\,u_t}\dd u}{t^{1/4} \bar L(\sqrt t)} \underset{t\rightarrow \infty} \longrightarrow \frac{1}{1+\gamma} \cdot \Phi(q)^{\frac{1}{2(1+\gamma)}-\frac{3}{4}} \int_0^{\infty}4u^2e^{-u^2}\dd u.
\]
It remains to show that $\int_0^{\infty} 2ue^{-u^2} I_{u,t,q} \ind{u\,>\,u_t}\dd u=o({t^{1/4} \bar L(\sqrt t)})$. In that aim, note that for $u>u_t$, one has $a_{u,t,q}+b_{u,t,q} \leq c_3(q)u^2$ for some constant $c_3(q)$ depending only on $q$. And also that $\bar V(x) \leq \kappa \sqrt x~$ for some $\kappa \in (0,\infty)$ and all $x$ large enough. This is clear by the regular variation property of $\bar V$ with an index less or equal to 1/2 (and more generally without the assumption of regular variation, as noticed in the proof of Proposition~\ref{prop:unifBound}).
Consequently, for $t$ large enough and $u>u_t$,
\begin{equation*}
I_{u,t,q} \leq \bar V \left(a_{u,t,q}+b_{u,t,q}\right) \leq \kappa \sqrt{c_3(q)} u.
\end{equation*}
This leads to
\begin{equation*}
\frac{\int_0^{\infty} 2ue^{-u^2} \ind{u\,>\,u_t}I_{u,t,q} \dd u}{t^{1/4} \bar L(\sqrt t)} =O \left(\frac{e^{-\frac{\sqrt{t \Phi(q)}}{8}}}{t^{1/4} \bar L(\sqrt t)} \right) \underset{t\rightarrow \infty} \longrightarrow 0
\end{equation*}
as required.
\end{proof}

\begin{thebibliography}{BGT87}
\bibitem[BGT87]{BGT} Bingham, Nick; Goldie, Charles; Teugels, Jozef Regular variation, Encyclopedia of Mathematics and Its Applications, 27, Cambridge University Press, 1987.
\bibitem[Ber98]{BLevy} Bertoin, Jean Lévy processes, Cambridge Tracts in Mathematics, 121, Cambridge University Press, 1998.
\bibitem[DR13]{DR13} Doney, Ron; Rivero, Victor Asymptotic behaviour of first passage time distributions for Lévy processes, Probab. Theory Relat. Fields, Volume 157 (2013) no. 1-2, pp. 1-45.
\bibitem[DR16]{DR16_err} Doney, Ron; Rivero, Victor Erratum to: “Asymptotic behaviour of first passage time distributions for Lévy processes”, Probab. Theory Relat. Fields, Volume 164 (2016) no. 3-4, pp. 1079-1083.
\bibitem[Kyp14]{KypFluctu} Andreas E. Kyprianou, Fluctuations of Lévy processes with applications. Introductory lectures, 2nd edition, Universitext, Springer,2014.
\end{thebibliography}
\end{document}
